Frontier coding agents pay per token, while Model Context Protocol (MCP) servers expose catalogs and repositories far
larger than any useful prompt. We study small local routers that explore such MCP surfaces and hand the frontier model
a compact, verifiable context packet. The routers are lexical, dense, feed-forward Transformer, shared-weight looped
Transformer, Tiny Recursive Model (TRM), TRM with certified lattice refinement, and HRM, from about 0.15 to 1.3M
parameters at recurrent depth 1 to 12. The paper reports three measurements:
\begin{itemize}
\item \textbf{Synthetic multi-server catalogs.} Routers with an exact-match token channel solve layered discovery tasks that BM25 and
hand-written rules mostly fail, using about half the tool calls. Looped and recursive cores match a four-layer
Transformer with a quarter of its non-embedding parameters, and deeper recurrence adds nothing measurable.
\item \textbf{SWE-bench Lite.} On \McpThreeInstances{} tasks from ten repositories unseen in training, a 2,048-token
packet that combines a BM25 path map with router-selected source snippets contains the gold files with recall
\BmComposedFileRecallTwoK{} and the gold functions with recall \BmComposedSymRecallTwoK{}. SWE-bench's BM25 whole-file
context reaches \BmFilesRecallTwentySeven{} file recall at \BmFilesTokensTwentySeven{} tokens.
\item \textbf{Cost.} Training a router fits in a few hundred megabytes of process memory, and inference adds megabytes on top of the
runtime.
\end{itemize}
We integrate the packet into Goose as an MCP extension and as a Rust app on the Goose Development Kit. With a local
0.5B model behind the app, the packet raises how often the model names the gold file. Every episode is stored with a
replayable decision tape.
